\chapter{Sparse Intervention Control}
\label{ch:sparse-control}

Every experiment so far has maximized one quantity, survival on held-out
chronics. This chapter changes the objective. A policy that keeps the grid
alive by reconfiguring substations at every opportunity is not what a control
room wants: interventions have an operational cost, and an agent that acts
constantly is neither auditable nor trustworthy. The question here is whether
the same decentralized agents can keep their survival while intervening rarely,
locally, and mainly when the grid state makes an intervention necessary.

The mechanisms are therefore evaluated on two axes at once, survival and how
often the agents leave the do-nothing action. A mechanism that improves
sparsity by collapsing survival has not solved the problem, and neither has one
that keeps survival by never becoming sparse.


\paragraph{Question.}
Can the agents keep high survival while behaving more like operators: acting
rarely, acting locally, acting mainly when the grid state makes an intervention
useful, and keeping the final policy decentralized?

\paragraph{Baseline and change.}
This block starts from the flat decentralized MAPPO topology-control actor. For
agent \(i\), the baseline policy is one categorical distribution over the full
local action space:
\[
  \pi_i(a_i \mid o_i),
  \qquad
  a_i \in \{0,1,\ldots,|\mathcal A_i|-1\}.
\]
Action \(0\) is the no-op action. During training, actions are sampled from the
policy, while deterministic evaluation uses the greedy action by default:
\[
  a_{i,t}^{eval}=\arg\max_a \pi_i(a\mid o_{i,t}).
\]
The sparse-control roadmap tests four ways to reduce unnecessary topology
changes: evaluation-time rho overrides, an intervention-gated actor, fixed
non-idle action penalties, and adaptive Lagrangian intervention budgets.

\paragraph{Motivation.}
The intervention gate is motivated by hierarchical topology-control work, where
the action can be decomposed into ``do nothing versus intervene'', then where to
act, then which local topology configuration to apply~\cite{manczak2023hrl,
vandersar2023hmarl}. In this thesis only the first level is adapted: each
regional agent independently decides whether to do nothing or to execute one of
its existing local topology actions. The fixed and adaptive sparse-control
objectives are closer to constrained and budgeted RL: non-idle topology
interventions are treated as a limited resource, following the general
Lagrangian view of constrained policy optimization and budgeted control
\cite{achiam2017cpo,stooke2020pid,carrara2019budgeted}. This is also
operationally meaningful because switching frequency and topology deviation are
real power-grid objectives, not just cosmetic metrics~\cite{lautenbacher2025morl}.

\section{Evaluation-Time Rho Heuristics}
\label{subsec:sparse-control-heuristics}

\paragraph{Implementation detail.}
The rho heuristic is an evaluation-only diagnostic. It does not change the
training objective and it does not prove that the policy learned sparse
behavior. The policy first proposes the deterministic action \(\bar a_{i,t}\),
then the evaluator may replace it with action \(0\).

The global version uses the pre-action maximum line loading
\[
  \rho_t^{global}=\max_{\ell\in\mathcal L}\rho_{\ell,t},
\]
and executes
\[
  a_{i,t}^{exec}=
  \begin{cases}
    0, & \rho_t^{global}<\rho_{safe},\\
    \bar a_{i,t}, & \text{otherwise,}
  \end{cases}
  \qquad \rho_{safe}=0.90.
\]
The local version is decentralized. Agent \(i\) computes the maximum rho on the
lines touching its regional substations,
\[
  \rho_{i,t}^{local}=\max_{\ell\in\mathcal L_i}\rho_{\ell,t},
\]
and only that agent is forced to no-op when
\(\rho_{i,t}^{local}<\rho_{safe}\). A boundary line can appear in the local
neighborhood of both adjacent agents, but no communication is required: both
agents can observe the line through their own physical neighborhood.

\paragraph{Interpretation.}
The heuristic is useful to ask what happens if safe-state interventions are
blocked, but it is not a learned method. Its action-0 rates must therefore be
reported as final executed evaluation actions, not as evidence that the policy
itself prefers action \(0\).

\section{Intervention-Gated Actor}
\label{subsec:sparse-control-gate}

\paragraph{Implementation detail.}
The gated actor changes the actor factorization. Instead of one categorical
head over all actions, each agent has a gate head and a non-idle action head:
\[
  \pi_i^{gate}(g_i\mid o_i),\qquad g_i\in\{0,1\},
\]
where \(0\) means do nothing and \(1\) means intervene, and
\[
  \pi_i^{nonidle}(b_i\mid o_i),\qquad
  b_i\in\{1,\ldots,|\mathcal A_i|-1\}.
\]
During training, the gate is sampled first. If \(g_{i,t}=0\), then
\(a_{i,t}=0\). If \(g_{i,t}=1\), the final action is sampled from the non-idle
head. The PPO log-probability of the executed action is therefore
\[
  \log P(a_i=0)=\log P(g_i=0),
\]
and, for \(a_i>0\),
\[
  \log P(a_i)=\log P(g_i=1)+\log P(b_i=a_i\mid g_i=1).
\]

Two deterministic decoders are used. The \texttt{final\_action\_map} decoder
selects the most likely final executed action:
\[
  P(a_i=0)=P(g_i=0),\qquad
  P(a_i=a>0)=P(g_i=1)P(b_i=a\mid g_i=1).
\]
The \texttt{hierarchical\_greedy} decoder first chooses the most likely gate
decision, then chooses the most likely non-idle action only if the gate selects
intervention.

The original coupled entropy term is
\[
  H_i=H(g_i)+P(g_i=1)H(b_i).
\]
This can accidentally reward intervention because larger \(P(g_i=1)\) unlocks
more non-idle entropy. The separated entropy variant is
\[
  H_i=\alpha_{gate}H(g_i)+\alpha_{nonidle}H(b_i),
\]
and is used in the gate-plus-budget diagnostic runs.

\paragraph{Interpretation.}
The gate is learned and decentralized, but it can restrict exploration more
strongly than a reward penalty. If it collapses early to no-op, the non-idle
head is rarely executed and receives little learning signal. For this reason,
the gate is treated as an important diagnostic rather than the main sparse
control contribution.

\section{Fixed Sparse-Action Penalty}
\label{subsec:sparse-control-fixed-penalty}

\paragraph{Implementation detail.}
The Sparse16 sweep adds a fixed reward penalty to the final executed action of
each agent:
\[
  r'_{i,t}=r_{i,t}
  -\lambda_{fixed}\mathbf{1}[a_{i,t}\neq0],
\]
with
\[
  \lambda_{fixed}\in\{0.0,0.001,0.003,0.01\}.
\]
The grid is
\[
  \text{actor}\in\{\text{flat},\text{gated}\},\qquad
  \lambda_{fixed}\in\{0.0,0.001,0.003,0.01\},\qquad
  \text{seed}\in\{0,1,2\},
\]
which gives \(2\times4\times3=24\) runs. All Sparse16 runs set
\texttt{safe\_intervention\_penalty = 0.0}, so the penalty is not rho-weighted.

\paragraph{Interpretation.}
Sparse16 changes only the training reward. It does not override actions during
evaluation and it does not hard-mask exploration during training. The policy can
still sample non-idle actions, but each such action must be useful enough to pay
its fixed cost. The key comparison is \texttt{sparse16\_flat\_p010} versus
\texttt{sparse16\_gated\_p010}, which asks whether the gate adds value once a
simple sparse objective already exists.

\section{Adaptive Intervention Budget}
\label{subsec:sparse-control-aib}

\paragraph{Implementation detail.}
The adaptive intervention budget turns sparse topology control into a local
constraint. For each agent,
\[
  c_i(s_t,a_{i,t})=
  \mathbf{1}[a_{i,t}\neq0]w_i(s_t),
\]
and the desired constraint is
\[
  \mathrm{E}_{\pi_i}\left[c_i(s_t,a_{i,t})\right]\le d.
\]
The implemented rollout reward removes the constant term that does not affect
the PPO action gradient:
\[
  r'_{i,t}=r_{i,t}-\lambda_i c_i(s_t,a_{i,t}).
\]
After each rollout, the local multiplier is updated by
\[
  \lambda_i\leftarrow
  \mathrm{clip}\left(
    \lambda_i+\eta(\hat C_i-d),
    0,
    \lambda_{\max}
  \right),
  \qquad
  \hat C_i=\frac{1}{T}\sum_{t=1}^{T}c_i(s_t,a_{i,t}).
\]
If the observed cost \(\hat C_i\) is above the budget \(d\), future
interventions become more expensive; if it is below the budget, the penalty
relaxes. The default settings are \texttt{intervention\_budget\_lr = 0.02},
\texttt{intervention\_budget\_init\_lambda = 0.0}, and
\texttt{intervention\_budget\_max\_lambda = 10.0}.

Three cost modes are used:
\begin{itemize}
  \item \textbf{Non-idle cost:}
  \(w_i(s_t)=1\), so
  \(c_i(s_t,a_{i,t})=\mathbf{1}[a_{i,t}\neq0]\). This is an adaptive version
  of a plain sparse-action penalty.
  \item \textbf{Global-safe cost:}
  \(w^{global}(s_t)=\sigma(k(\rho_{safe}-\rho_t^{global}))\). This penalizes
  non-idle actions mostly when the whole grid is safe.
  \item \textbf{Local-safe cost:}
  \(w_i^{local}(s_t)=\sigma(k(\rho_{safe}-\rho_{i,t}^{local}))\). This is the
  decentralized version: safe local states make interventions expensive, while
  stressed local states make them cheap.
\end{itemize}
In the AIB setting, \(\rho_{safe}=0.90\) is not a hard action override. It is
the center of a smooth sigmoid safety weight.

\paragraph{AIB experiment grid.}
The first AIB folder, \texttt{configs/adaptive\_intervention\_budget\_7},
contains the main local-budget ablations:
\begin{itemize}
  \item \texttt{aib\_00\_flat\_local\_t020}: flat actor, local-safe cost,
  \(d=0.20\);
  \item \texttt{aib\_01\_flat\_local\_t010}: stricter local-safe budget,
  \(d=0.10\);
  \item \texttt{aib\_02\_flat\_local\_t035}: looser local-safe budget,
  \(d=0.35\);
  \item \texttt{aib\_03\_gate\_hgreedy\_sep\_local\_t020}: gated actor,
  hierarchical-greedy evaluation, separated entropy, local-safe cost,
  \(d=0.20\);
  \item \texttt{aib\_04\_flat\_nonidle\_t020}: flat actor with adaptive plain
  non-idle cost, \(d=0.20\).
\end{itemize}
The comparison \texttt{aib\_00\_flat\_local\_t020} versus
\texttt{aib\_04\_flat\_nonidle\_t020} isolates the value of local rho weighting.

The mechanism-ablation folder
\texttt{configs/adaptive\_intervention\_budget\_mechanism\_15} contains 15
additional flat-actor runs. It separates fixed versus adaptive coefficients,
plain non-idle versus state-dependent costs, global versus local rho weighting,
and rho-threshold sensitivity.

\begin{table}[htbp]
  \centering
  \small
  \caption{Sparse-control mechanisms and where they affect the policy.}
  \label{tab:sparse-control-mechanisms}
  \setlength{\tabcolsep}{3pt}
  \begin{tabular}{L{0.18\textwidth}L{0.30\textwidth}L{0.24\textwidth}L{0.18\textwidth}}
    \toprule
    Method & Training effect & Evaluation effect & Main risk \\
    \midrule
    Baseline & Samples from the flat policy & Greedy argmax by default &
    Standard entropy-controlled exploration \\
    Rho heuristic & No training change & Hard no-op override in safe states &
    Evaluation is manually changed \\
    Intervention gate & Changes action factorization and sampling &
    Final-action MAP or hierarchical greedy & Gate collapse can starve non-idle exploration \\
    Sparse16 & Fixed reward penalty on non-idle actions & No action override &
    Penalty coefficient must be tuned \\
    AIB non-idle & Adaptive reward penalty on non-idle actions & No action override &
    Multiplier can grow if target is too strict \\
    AIB local-safe & Adaptive state-dependent reward penalty & No action override &
    Depends on the local rho safety-weight design \\
    \bottomrule
  \end{tabular}
\end{table}

\begin{table}[htbp]
  \centering
  \small
  \caption{Main sparse-control comparisons to report.}
  \label{tab:sparse-control-main-comparisons}
  \setlength{\tabcolsep}{4pt}
  \begin{tabular}{L{0.45\textwidth}L{0.45\textwidth}}
    \toprule
    Question & Comparison \\
    \midrule
    Does a fixed sparse penalty already work? &
    Baseline versus \texttt{sparse16\_flat\_p010} \\
    Does the gate help beyond a sparse objective? &
    \texttt{sparse16\_flat\_p010} versus \texttt{sparse16\_gated\_p010} \\
    Fixed versus adaptive coefficient? &
    \texttt{sparse16\_flat\_p010} or \texttt{aibm\_00\_fixed\_nonidle\_p006}
    versus \texttt{aib\_04\_flat\_nonidle\_t020} \\
    Does local-safe weighting matter? &
    \texttt{aib\_00\_flat\_local\_t020} versus
    \texttt{aib\_04\_flat\_nonidle\_t020} \\
    Does global-safe weighting matter? &
    \texttt{aibm\_02\_adaptive\_global\_t020} versus
    \texttt{aib\_00\_flat\_local\_t020} \\
    Is \(\rho_{safe}=0.90\) special? &
    \texttt{aibm\_03} with \(0.85\), \texttt{aib\_00} with \(0.90\),
    and \texttt{aibm\_04} with \(0.95\) \\
    Is the heuristic only an evaluation trick? &
    Rho-heuristic runs versus learned sparse/AIB runs \\
    \bottomrule
  \end{tabular}
\end{table}

\paragraph{Full-test checkpoint summaries.}
The following tables report the average full-test episodic survival over the
three evaluated seeds. The survival values are computed directly from the JSON
files in \texttt{Topology\_Task/outputs/full\_test\_eval}. These JSON files also
record where the compact action-distribution artifacts were written. The
action-0 fractions are computed from the corresponding local
\texttt{Topology\_Task/outputs/full\_test\_eval\_actions/*/action\_summary.json}
files and averaged over seeds. The all-agent action-0 column is the arithmetic
mean of the three per-agent action-0 fractions. The do-nothing row is a
reference policy that always selects action 0; its survival is computed from
\texttt{Topology\_Task/outputs/do\_nothing\_eval/bus14\_test\_20}. Remaining
``--'' entries indicate that the compact action summaries for those full-test
jobs are not available locally.

\begin{table}[htbp]
  \centering
  \small
  \caption{Full-test summary for evaluation-time rho heuristics.}
  \label{tab:a0-heuristic-full-test}
  \setlength{\tabcolsep}{4pt}
  \resizebox{\textwidth}{!}{%
  \begin{tabular}{L{0.34\textwidth}rrrrr}
    \toprule
    Run type & Avg. episodic survival (\%) & Avg. action-0 all agents &
    Avg. action-0 agent 0 & Avg. action-0 agent 1 &
    Avg. action-0 agent 2 \\
    \midrule
    Do-nothing policy & 20.25 & 1.000 & 1.000 & 1.000 & 1.000 \\
    Baseline flat policy & 93.67 & 0.434 & 0.554 & 0.385 & 0.363 \\
    Global rho heuristic, \(\rho_{safe}=0.90\) & 98.36 & 0.936 & 0.933 & 0.936 & 0.940 \\
    Local rho heuristic, \(\rho_{safe}=0.90\) & 97.07 & 0.968 & 0.992 & 0.985 & 0.928 \\
    \bottomrule
  \end{tabular}
  }
\end{table}

\begin{table}[htbp]
  \centering
  \small
  \caption{Full-test summary for intervention-gated actors.}
  \label{tab:a0-gate-full-test}
  \setlength{\tabcolsep}{4pt}
  \resizebox{\textwidth}{!}{%
  \begin{tabular}{L{0.34\textwidth}rrrrr}
    \toprule
    Run type & Avg. episodic survival (\%) & Avg. action-0 all agents &
    Avg. action-0 agent 0 & Avg. action-0 agent 1 &
    Avg. action-0 agent 2 \\
    \midrule
    Do-nothing policy & 20.25 & 1.000 & 1.000 & 1.000 & 1.000 \\
    Baseline flat policy & 93.67 & 0.434 & 0.554 & 0.385 & 0.363 \\
    Intervention gate, final-action MAP & 89.82 & 0.802 & 0.826 & 0.889 & 0.690 \\
    Intervention gate, hierarchical greedy & 93.66 & 0.395 & 0.066 & 0.844 & 0.276 \\
    \bottomrule
  \end{tabular}
  }
\end{table}

\begin{table}[htbp]
  \centering
  \small
  \caption{Full-test summary for Sparse16 flat-policy runs.}
  \label{tab:a0-sparse16-flat-full-test}
  \setlength{\tabcolsep}{4pt}
  \resizebox{\textwidth}{!}{%
  \begin{tabular}{L{0.34\textwidth}rrrrr}
    \toprule
    Run type & Avg. episodic survival (\%) & Avg. action-0 all agents &
    Avg. action-0 agent 0 & Avg. action-0 agent 1 &
    Avg. action-0 agent 2 \\
    \midrule
    Do-nothing policy & 20.25 & 1.000 & 1.000 & 1.000 & 1.000 \\
    Baseline flat policy & 93.67 & 0.434 & 0.554 & 0.385 & 0.363 \\
    Flat Sparse16, \(\lambda_{fixed}=0.000\) & 98.23 & 0.494 & 0.524 & 0.463 & 0.496 \\
    Flat Sparse16, \(\lambda_{fixed}=0.001\) & 96.75 & 0.486 & 0.642 & 0.401 & 0.416 \\
    Flat Sparse16, \(\lambda_{fixed}=0.003\) & 93.97 & 0.575 & 0.691 & 0.346 & 0.688 \\
    Flat Sparse16, \(\lambda_{fixed}=0.010\) & 97.44 & 0.905 & 0.907 & 0.954 & 0.853 \\
    \bottomrule
  \end{tabular}
  }
\end{table}

\begin{table}[htbp]
  \centering
  \small
  \caption{Full-test summary for Sparse16 gated-policy runs.}
  \label{tab:a0-sparse16-gated-full-test}
  \setlength{\tabcolsep}{4pt}
  \resizebox{\textwidth}{!}{%
  \begin{tabular}{L{0.34\textwidth}rrrrr}
    \toprule
    Run type & Avg. episodic survival (\%) & Avg. action-0 all agents &
    Avg. action-0 agent 0 & Avg. action-0 agent 1 &
    Avg. action-0 agent 2 \\
    \midrule
    Do-nothing policy & 20.25 & 1.000 & 1.000 & 1.000 & 1.000 \\
    Baseline flat policy & 93.67 & 0.434 & 0.554 & 0.385 & 0.363 \\
    Gated Sparse16, \(\lambda_{fixed}=0.000\) & 78.51 & 0.787 & 0.747 & 0.949 & 0.665 \\
    Gated Sparse16, \(\lambda_{fixed}=0.001\) & 73.82 & 0.862 & 0.956 & 0.832 & 0.799 \\
    Gated Sparse16, \(\lambda_{fixed}=0.003\) & 94.59 & 0.887 & 0.934 & 0.929 & 0.799 \\
    Gated Sparse16, \(\lambda_{fixed}=0.010\) & 85.69 & 0.974 & 0.990 & 0.984 & 0.947 \\
    \bottomrule
  \end{tabular}
  }
\end{table}

\begin{table}[htbp]
  \centering
  \small
  \caption{Full-test summary for adaptive intervention budget runs.}
  \label{tab:a0-aib-full-test}
  \setlength{\tabcolsep}{4pt}
  \resizebox{\textwidth}{!}{%
  \begin{tabular}{L{0.34\textwidth}rrrrr}
    \toprule
    Run type & Avg. episodic survival (\%) & Avg. action-0 all agents &
    Avg. action-0 agent 0 & Avg. action-0 agent 1 &
    Avg. action-0 agent 2 \\
    \midrule
    Do-nothing policy & 20.25 & 1.000 & 1.000 & 1.000 & 1.000 \\
    Baseline flat policy & 93.67 & 0.434 & 0.554 & 0.385 & 0.363 \\
    Flat local-safe AIB, \(d=0.20\) & 98.31 & 0.978 & 0.980 & 0.989 & 0.966 \\
    Flat local-safe AIB, \(d=0.10\) & 98.39 & 0.979 & 0.992 & 0.995 & 0.949 \\
    Flat local-safe AIB, \(d=0.35\) & 96.61 & 0.934 & 0.959 & 0.947 & 0.896 \\
    Gated local-safe AIB, \(d=0.20\), hierarchical greedy &
    33.87 & 0.999 & 1.000 & 1.000 & 0.998 \\
    Flat non-idle AIB, \(d=0.20\) & 95.08 & 0.986 & 0.997 & 0.992 & 0.970 \\
    \bottomrule
  \end{tabular}
  }
\end{table}

\begin{table}[htbp]
  \centering
  \small
  \caption{Full-test summary for the teacher-student distillation runs.}
  \label{tab:teacher-student-full-test}
  \setlength{\tabcolsep}{4pt}
  \resizebox{\textwidth}{!}{%
  \begin{tabular}{L{0.34\textwidth}rrrrr}
    \toprule
    Run type & Avg. episodic survival (\%) & Avg. action-0 all agents &
    Avg. action-0 agent 0 & Avg. action-0 agent 1 &
    Avg. action-0 agent 2 \\
    \midrule
    Do-nothing policy & 20.25 & 1.000 & 1.000 & 1.000 & 1.000 \\
    Teacher: MAPPO with local rho heuristic & 97.07 & 0.968 & 0.992 & 0.985 & 0.928 \\
    Student BC, balanced non-idle 0.50, weight 5, aux 0.50 &
    92.14 & 0.939 & 0.974 & 0.925 & 0.919 \\
    Student BC, balanced non-idle 0.20, weight 3, aux 0.25 &
    90.76 & 0.957 & 0.982 & 0.948 & 0.940 \\
    \bottomrule
  \end{tabular}
  }
\end{table}

\paragraph{Metrics to report.}
The main performance metric remains episodic survival. It should be paired with
sparsity and physical-diagnostic metrics: per-agent action-0 fractions,
per-agent non-idle rates, the distribution of how many agents act
simultaneously, intervention-budget multipliers, budget costs and violations,
the rate of interventions in costly safe states, pre- and post-action max rho,
and topology-distance changes.

\paragraph{Interpretation.}
The current roadmap interpretation is that the gate is not the strongest
contribution by itself. Sparse reward shaping already improves action-0 usage,
the adaptive budget makes the sparse penalty self-tuning, and the local-safe
cost makes the penalty more physically meaningful. In particular,
\texttt{sparse16\_flat\_p010} is a strong fixed-penalty baseline, while the
most principled learned sparse-control candidate is
\texttt{aib\_00\_flat\_local\_t020}: it keeps the original decentralized actor,
does not rely on an evaluation-time override, learns the intervention penalty
adaptively, and uses a local state-dependent safety weight.


\section{Preliminary WCCI Stress Test}
\label{sec:sparse-wcci-stress-test}

As a preliminary scale test, representative mechanisms were retrained from
scratch with an MLP actor on the maintenance-enabled WCCI grid; no \texttt{bus14}
weights were transferred. All runs use the same reduced action space, three
seeds, 60M training steps, and learning-rate schedule. The reported values are
the final deterministic test evaluation, averaged over seeds and ten held-out
episodes per seed; they are not full-test estimates.

\begin{table}[htbp]
  \centering
  \small
  \caption{Preliminary WCCI stress test with scheduled maintenance at the final
  training-time test evaluation.}
  \label{tab:sparse-wcci-stress-test}
  \setlength{\tabcolsep}{7pt}
  \begin{tabular}{L{0.50\textwidth}rr}
    \toprule
    Method & Survival (\%) & Action-0 fraction \\
    \midrule
    Matched MLP baseline & $23.75\pm7.70$ & 0.168 \\
    Global rho heuristic, $\rho_{safe}=0.90$ & $28.35\pm11.06$ & 0.974 \\
    Local rho heuristic, $\rho_{safe}=0.90$ & $9.46\pm3.54$ & 0.995 \\
    Intervention gate, final-action MAP & $25.92\pm10.69$ & 0.400 \\
    Flat Sparse16, $\lambda_{fixed}=0.010$ & $27.00\pm12.13$ & 0.489 \\
    Flat local-safe AIB, $d=0.20$ & $21.65\pm2.21$ & 0.927 \\
    \bottomrule
  \end{tabular}
\end{table}

The mechanisms change intervention frequency, but none reliably improves
survival over the weak MLP baseline: the differences are small relative to the
seed variation. The local rho heuristic and AIB are highly sparse but remain
weak controllers, suggesting that sparsity transfers more readily than control
quality. This failed pilot changes grid size, action space, and maintenance at
once, so it cannot identify the cause. Chapter~\ref{ch:wcci-transfer} isolates
these factors before studying cross-grid graph transfer.
